% ======================================================================
\section{Discussion}
\label{sec:discussion}

\subsection{What the evidence supports}
The three contributions were stated as claims that an experiment could falsify, and the
results section reports the outcome of each. The phase-locked tier is a component, not a
model: it is worth \nv{abl.patch.no_pl.mean_dwape_pct} and
\nv{abl.basis.no_pl.mean_dwape_pct} in WAPE to two residual backbones that share nothing but
the interface, averaged over the reference cells, and \nv{abl.patch.OPSD-DE.no_pl.H24.dwape_pct}
and \nv{abl.basis.OPSD-DE.no_pl.H24.dwape_pct} on the German system at one day, and it moves both backbones in the same direction in every cell but two, where its effect on
one backbone is within a point of zero. On the eight-system benchmark
the model built around it is retained by the Model Confidence Set in
\nv{mcs.grid.PLANET.n_in} of \nv{mcs.grid.n_cells} cells and has the second mean rank of
\nv{grid.n_models} methods; it is significantly worse than a compared method in
\nv{dm.total.losses} of \nv{dm.n_tests} tests, \nv{dm.TiDE.losses} of them against TiDE, which
is the strongest form of accuracy claim the data support and the only one we make. The
mechanism claim held where the mechanism says it should, against the learned-calendar
control: per-window estimation pays when the model is moved to a system it has never seen,
and with one quarter to one half of the history, where it keeps PLANET ahead of TiDE on
Germany and carries it past TiDE on Spain with a quarter of the history but not on France, and the size of both effects is
consistent with the explanation, that a calendar map must be learned from examples of every
day-type while a calendar basis need not. The protocol claim is supported by its own data:
handing a fixed architecture the calendar and the weather forecast changes its error by
\nv{info.grid.R2_vs_R0.mean_dwape_pct}, more than the gap between the proposed model and the
median architecture given the same inputs, so that comparisons which do not control the
information set are comparisons of interfaces.

\subsection{What the evidence does not support}
Four findings go against hypotheses that motivated the design, and we report them as results
rather than drop them. The first is that a closed-form calendar fit is more accurate than a
learned calendar map when training data are plentiful: on the full training split
PatchTST+Cal, which is handed the same design matrix as tokens, is statistically
indistinguishable from PLANET in \nv{dm.CalTST.ties} of \nv{dm.CalTST.n_cells} cells, and
N-BEATS, which has no calendar input beyond the shared head, is indistinguishable in
\nv{dm.NBEATS.ties}. A reader who has four years of clean history and no intention of moving
the model may reasonably prefer whichever of the three is simplest to maintain. The second is
that gradient-free adaptation alone delivers robustness to regime shift: under temporal drift
every method degrades by a comparable amount once it has half the history, because what
drifts is the residual, not the seasonal state that the solve re-estimates. A model deployed
across a structural break still needs to be retrained or to carry an adaptive residual, and
the tier does not remove that need. The third is that the two defining choices of the basis,
locking it to the calendar rather than the window and evaluating it on local civil time, are
where its value lies: the first is worth \nv{abl.patch.relative.mean_dwape_pct} and
\nv{abl.basis.relative.mean_dwape_pct} on average and the second nothing measurable, so the
tier's value is in the per-window solve of a calendar basis as such, not in those two
refinements. The fourth, and the one a reader will weigh most, is that the proposed model is
not the most accurate forecaster in its own benchmark. TiDE is significantly ahead of it in
\nv{dm.TiDE.losses} cells and behind in none, with a mean rank of \nv{grid.TiDE.meanrank_wape}
against \nv{grid.PLANET.meanrank_wape}; the next paragraph says why, and the reason is
specific enough to be useful.

\paragraph{Where the calendar tier pays, and where the weather pathway does}
\tabv{tab_weather_v2}
The natural explanation of the headline table is a contrast between
calendar-driven and weather-driven systems: German load is a calendar signal (daily mean load
correlates with daily temperature at \nv{data.OPSD-DE.wx_corr}, winter exceeds summer by a
factor \nv{data.OPSD-DE.wx_winter_summer}, and the weather block changes the four-day error of
the gradient-boosted trees by \nv{info.OPSD-DE.GBRT.H96.R2_vs_R1.dwape_pct}), French load a
weather signal (electric heating; correlation \nv{data.OPSD-FR.wx_corr}, factor
\nv{data.OPSD-FR.wx_winter_summer}, weather gain \nv{info.OPSD-FR.GBRT.H96.R2_vs_R1.dwape_pct}
for the same trees); PLANET leads every model at every horizon on the first and trails TiDE on
the second. Table~\ref{tab:weather-v2} puts that explanation beside all eight systems, ordered
by the share of daily load that temperature explains, and it does not hold as a rule: the Dutch
and Polish systems are as weakly weather-dependent as the German one, the Italian one less so,
and TiDE leads PLANET on all three. What the
table does support is narrower: the tier is worth most on Germany
(\nv{abl.patch.OPSD-DE.no_pl.H96.dwape_pct} at four days on the same backbone against
\nv{abl.patch.OPSD-FR.no_pl.H96.dwape_pct} on France), PLANET's standing against the
learned-calendar control is favourable on six of the eight systems, within a point on Great Britain and unfavourable
on Poland (last column), and its deficit to TiDE is a property of the pair, not of the weather. The pathway control of Section~\ref{sec:mech}
puts a number on that pair: routing TiDE's covariates through the shared head instead of its
per-step encoder changes TiDE's WAPE by \nv{path.grid.TiDEh.mean_dwape_vs_TiDE_pct} on
average, and once both models read the covariates through the same head the remaining
difference to PLANET is \nv{path.grid.TiDEh.mean_dwape_vs_PLANET_pct}. PLANET-X, which gives
the tier a native pathway of its own, was pre-registered as the direct test of the pathway
explanation and failed it: its four-day WAPE relative to PLANET is
\nv{ref.OPSD-FR.H96.vs.PLANETX.gain_wape_pct} on France and
\nv{ref.OPSD-DE.H96.vs.PLANETX.gain_wape_pct} on Germany (positive: PLANET better), a marginal
gain on the weather-driven system and a clear loss on the calendar-driven one. A dense encoding
of the covariate block concatenated to the residual input is therefore not the pathway that
TiDE's per-step temporal decoder is. We did not iterate on further variants, since a variant
chosen after seeing the test grid would be worth nothing; the per-step pathway is the obvious
next design, and we leave it to be tested under the same pre-registration.

The paper also does not show that the operator can be replaced. The learned models beat the
published day-ahead forecast on \nv{da.grid.any.n_beats_tso_unbiased} of the
\nv{da.grid.n_unbiased} systems whose published series is unbiased, PLANET on
\nv{da.grid.PLANET.n_beats_tso_unbiased} and TiDE on \nv{da.grid.TiDE.n_beats_tso_unbiased},
under a like-for-like cut-off, but they hold degraded reanalysis rather than an archived
operational weather forecast, and the operator's forecast is a product of a process, including
manual correction, that a single model does not reproduce. The like-for-like result shows that
the protocol is fair, and that the accuracy an operator obtains is reachable by a model trained
under it; it does not show more than that.

\subsection{Threats to validity}
\paragraph{The weather proxy} We did not have access to an archive of the operational
day-ahead weather forecasts issued for the eight systems over 2015--2020 at hourly resolution
(archived global forecasts exist at coarser time steps and would be a natural extension), so
the future weather block is reanalysis degraded by a lead-dependent error model whose
parameters are set at the level of published day-ahead skill: for temperature the error
model reaches about two kelvin at 72\,h, the error standard deviation that ECMWF reports for
its 60--72\,h two-metre temperature forecasts against European surface stations
\cite{haiden2022ecmwf}, and for wind and solar generation it uses the typical day-ahead error
of national aggregates, several percent of installed capacity
\cite{hodge2012wind,lorenz2011pv}, whose real distributions are heavier-tailed than the
Gaussian assumed here. The error model is unbiased and Gaussian, which is kinder than a real forecast in every respect that matters here: no systematic
errors, no missed fronts, and errors that are independent across hours and variables where a
real forecast's persist, a difference the correlated-noise check of
Section~\ref{sec:sensitivity} prices at \nv{exo.ar.grid.TiDE.mean_dwape_pct} and
\nv{exo.ar.grid.PLANET.mean_dwape_pct}. Every method receives
the same proxy, so the comparison between methods is unaffected; the absolute level of R2
error, and the comparison with the operator, should be read with this caveat.

\paragraph{The training budget} Every model is trained with a cosine schedule to early stopping
on validation loss under a cap of thirty epochs; \nv{proto.epochs.frac_capped} of the
\nv{proto.epochs.n_runs} headline runs ran to the cap, which a cosine schedule makes expected
rather than alarming, since the validation loss keeps improving slowly until the schedule ends.
The relevant question is whether the budget shapes the ordering. We re-trained the ten
headline models on the German system with three times the cap (ninety epochs, patience ten,
three seeds): their WAPE changed by \nv{budget.OPSD-DE.mean_dwape_pct} on average (between
\nv{budget.OPSD-DE.min_dwape_pct} and \nv{budget.OPSD-DE.max_dwape_pct} over
\nv{budget.OPSD-DE.n_cells} cells), the Spearman correlation between the rankings under the
two budgets is \nv{budget.OPSD-DE.mean_spearman} averaged over horizons, and the best model
changed in \nv{budget.OPSD-DE.n_best_changed} of \nv{budget.OPSD-DE.n_horizons} horizons.
The longer budget does not help: most cells are slightly less accurate under it, so the cap is
acting as a regulariser rather than a constraint, the ordering is preserved, and the one change
at the top of the German table runs against the proposed model, since at two weeks \nv{budget.OPSD-DE.H336.best_large} edges past
\nv{budget.OPSD-DE.H336.best_small} once both train three times longer, by a small margin
(\nv{budget.OPSD-DE.TiDE.H336.wape} against \nv{budget.OPSD-DE.PLANET.H336.wape} WAPE). The French check (Section~\ref{sec:sensitivity}) reads the same way:
the best model changes at one horizon there too, and again away from the proposed model, with
\nv{budget.OPSD-FR.H24.best_large} overtaking \nv{budget.OPSD-FR.H24.best_small} at one day.
The budget therefore does not shape the ordering in the proposed model's favour.

\paragraph{The test span} The test period, August 2019 to September 2020, contains one
year-end holiday period, which favours a method with a holiday basis, and the spring-2020
lockdowns, which favour nothing that was trained before them. Both are in the test span of
every method and the per-window error files released with the paper allow either to be
excluded. We chose not to exclude them: an operational forecaster is scored on the year it
gets.

\paragraph{Foundation-model contamination} The pre-training corpora of the two zero-shot
models include public electricity series and may include the systems studied here. We have
reported them as references and drawn no conclusion from their absolute accuracy.

\paragraph{Scope} Eight European national systems with a shared regulatory environment and
similar load composition are not the world. The mechanism does not depend on Europe, but the
size of the holiday effect, the weather response and the value of the weekly cycle do, and a
system with a different working week or a different climate is a different experiment.

\subsection{Implications for practice}
For an operator or a forecasting team the results reduce to five points. The information set
matters more than the architecture: any competent model given calendar, holidays and the
weather forecast will outperform a better model without them, and the shared head of
Section~\ref{sec:head} is a cheap way to give any model that interface. Among models with the same inputs, a model whose main encoder reads the covariates, TiDE in
this study, is the most accurate point and probabilistic forecaster on most systems and among
the cheapest to train, so it, or its equal-weight mean with PLANET, is the default; the
closed-form calendar tier is the choice where the history is short or the model must be moved
between systems, and it costs \nv{cost.PLANET_vs_backbone.infer_pct} in inference against its
backbone, a small linear solve per window. The two are complementary: their equal-weight mean is more accurate than
either alone in nearly every cell of the grid (Section~\ref{sec:sensitivity}), and under the
operator protocol the combined ensemble improves on PLANET's ensemble on every system and on
TiDE's on most, tying it on the rest (Table~\ref{tab:dayahead-v2}), so a team that can afford
to run both should average them. A
model that must survive a structural break needs retraining or an adaptive residual, whatever
its seasonal representation. And calibrated intervals are a property of the calibrator, not the
forecaster: conformalised quantile regression followed by adaptive conformal inference brought
every method here to nominal coverage, after which the choice of forecaster is a choice of
sharpness.
